\documentclass[11pt]{amsart}
\usepackage[margin=28mm]{geometry}
\usepackage{amssymb,tikz,url}
\usepackage[hidelinks]{hyperref}
\newtheorem{theorem}{Theorem}
\newtheorem{proposition}[theorem]{Proposition}
\theoremstyle{remark}
\newtheorem*{remark}{Remark}

\title{The $697\times611$ rectangle needs exactly fourteen squares}
\author{Vamshi Jandhyala}
\date{Preprint, 9 October 2026. Not peer reviewed.}

\begin{document}
\begin{abstract}
Let $f(m,n)$ be the least number of integer-sided squares that tile an $m\times n$ rectangle. A 2017
answer on MathOverflow named $697\times611$ as the one rectangle with sides at most $760$ that might
violate the conjecture that scaling a rectangle never lowers $f$: it had a known tiling by $17$
squares, while its double $1394\times1222$ had one by $16$. We show that $f(697,611)=14$. The tiling is
explicit; the lower bound comes from enumerating every squared rectangle with at most $13$ squares
through the electrical networks of Brooks, Smith, Stone and Tutte, an enumeration that reproduces all
known values of $f$ for sides up to $388$.
\end{abstract}
\maketitle

\section{The question}

Cut a rectangle with integer sides into squares with integer sides. How few squares are needed? Write
$f(m,n)$ for the answer for an $m\times n$ rectangle. The values for all $m\le n\le 388$ are known and
tabulated in the OEIS~\cite{oeis}.

Scaling a tiling of an $m\times n$ rectangle by $t$ tiles the $tm\times tn$ rectangle with the same
number of squares, so $f(tm,tn)\le f(m,n)$. It is natural to expect equality, and this is sometimes
called the minimal squaring conjecture. In 2017, answering a MathOverflow question about the values of
$f$~\cite{mo}, Ed Pegg Jr reported that among rectangles with sides at most $760$ one possible
counterexample remained: $697\times 611$, for which the best tiling known used $17$ squares, while
$1394\times1222$ could be tiled with $16$~\cite{pegg}. If $17$ were the minimum, doubling would save a
square.

\begin{theorem}\label{thm:main}
$f(697,611)=14$. In particular $697\times611$ is not a counterexample to the minimal squaring
conjecture.
\end{theorem}

\section{Fourteen squares suffice}

Figure~\ref{fig:tiling} shows the tiling. Its squares have sides
\[
371,\ 326,\ 285,\ 240,\ 172,\ 68,\ 41,\ 37,\ 35,\ 35,\ 34,\ 34,\ 33,\ 4,
\]
and their areas add up to $425\,867=697\cdot611$. With the origin at the bottom left corner, the
squares sit at the positions listed in \path{paper/make_fig.py}, which checks cell by cell that they
cover the rectangle without overlapping before it draws the figure. In Bouwkamp code, reading the
squares row by row from the top, the tiling is
\[
(240,172,285)\,(68,34,35,35)\,(34)\,(33,37)\,(371,4)\,(41)\,(326).
\]

\begin{figure}[ht]
\centering
\begin{tikzpicture}
\filldraw[fill=black!6,draw=black,line width=0.5pt] (0.0000,0.0000) rectangle (5.3228,5.3228);
\node[font=\small] at (2.6614,2.6614) {371};
\filldraw[fill=black!6,draw=black,line width=0.5pt] (5.3228,0.0000) rectangle (10.0000,4.6772);
\node[font=\small] at (7.6614,2.3386) {326};
\filldraw[fill=black!6,draw=black,line width=0.5pt] (5.3228,4.6772) rectangle (5.9110,5.2654);
\node[font=\tiny] at (5.6169,4.9713) {41};
\filldraw[fill=black!6,draw=black,line width=0.5pt] (5.9110,4.6772) rectangle (10.0000,8.7661);
\node[font=\small] at (7.9555,6.7217) {285};
\filldraw[fill=black!6,draw=black,line width=0.5pt] (5.3228,5.2654) rectangle (5.3802,5.3228);
\filldraw[fill=black!6,draw=black,line width=0.5pt] (5.3802,5.2654) rectangle (5.9110,5.7963);
\node[font=\tiny] at (5.6456,5.5308) {37};
\filldraw[fill=black!6,draw=black,line width=0.5pt] (0.0000,5.3228) rectangle (3.4433,8.7661);
\node[font=\small] at (1.7217,7.0445) {240};
\filldraw[fill=black!6,draw=black,line width=0.5pt] (3.4433,5.3228) rectangle (4.4189,6.2984);
\node[font=\tiny] at (3.9311,5.8106) {68};
\filldraw[fill=black!6,draw=black,line width=0.5pt] (4.4189,5.3228) rectangle (4.9067,5.8106);
\node[font=\tiny] at (4.6628,5.5667) {34};
\filldraw[fill=black!6,draw=black,line width=0.5pt] (4.9067,5.3228) rectangle (5.3802,5.7963);
\node[font=\tiny] at (5.1435,5.5595) {33};
\filldraw[fill=black!6,draw=black,line width=0.5pt] (4.9067,5.7963) rectangle (5.4089,6.2984);
\node[font=\tiny] at (5.1578,6.0473) {35};
\filldraw[fill=black!6,draw=black,line width=0.5pt] (5.4089,5.7963) rectangle (5.9110,6.2984);
\node[font=\tiny] at (5.6600,6.0473) {35};
\filldraw[fill=black!6,draw=black,line width=0.5pt] (4.4189,5.8106) rectangle (4.9067,6.2984);
\node[font=\tiny] at (4.6628,6.0545) {34};
\filldraw[fill=black!6,draw=black,line width=0.5pt] (3.4433,6.2984) rectangle (5.9110,8.7661);
\node[font=\small] at (4.6772,7.5323) {172};
\draw[line width=1.2pt] (0,0) rectangle (10.0000,8.7661);
\end{tikzpicture}

\caption{A tiling of the $697\times611$ rectangle by fourteen squares. Two pairs of equal squares,
$35,35$ and $34,34$, sit side by side.}\label{fig:tiling}
\end{figure}

The tiling is \emph{compound}: it contains a smaller rectangle tiled by squares, for instance the
$70\times35$ rectangle formed by the two squares of side $35$, and it uses equal squares. Searches of
catalogues of simple perfect squared rectangles, which contain neither feature, cannot find it.

\section{Thirteen squares do not}

The lower bound uses the correspondence of Brooks, Smith, Stone and Tutte~\cite{bsst} between squared
rectangles and electrical networks. Given a rectangle tiled by $k$ squares, take a vertex for each
maximal horizontal segment of the tiling, including the top and bottom sides, and an edge for each
square, joining the segment its top lies on to the segment its bottom lies on. Add one more edge, the
\emph{pole}, from the top side to the bottom side. The result is a $2$-connected plane multigraph
with $k+1$ edges. Put a unit resistor on each square's edge and a battery on the pole: the current
through each edge is the side of its square, and the potential difference across the battery is the
height of the rectangle. Conversely, a $2$-connected plane multigraph with a chosen pole edge, in which
no current vanishes, yields a squared rectangle in this way.

So every tiling of a rectangle by $k$ squares appears among the networks with $k+1$ edges, and to rule
out a tiling it suffices to examine all of them. The $2$-connected plane multigraphs with $e$ edges
are exactly one colour class of the simple quadrangulations of the sphere with $e$ faces; the other
class gives the dual networks, which produce the same tilings turned through a right angle. These
quadrangulations have $e+2$ vertices and are generated by the program \emph{plantri} of Brinkmann
and McKay~\cite{plantri}.

\begin{proposition}[computer-certified]\label{prop:lower}
No network with at most $14$ edges produces a tiling whose rectangle is similar to $697\times611$.
\end{proposition}

\begin{proof}[Method]
The program \path{verify/sqrect.c} reads every quadrangulation with $5$ to $16$ vertices from
plantri, takes each edge in turn as the pole, and solves for the currents. The currents are computed in
floating point, scaled by the number of spanning trees and rounded, and then verified exactly in integer
arithmetic: with $L$ the Laplacian and $D$ the number of spanning trees, the rounded potentials $N$
must satisfy $LN=Db$. No run has failed this check. Each tiling is reduced by the greatest common
divisor of its square sides and its shape recorded. Since $697=17\cdot41$ and $611=13\cdot47$ are
coprime, a tiling of $697\times611$ by $k\le13$ squares would be recorded with exactly this shape. None
is: the shape first appears with $14$ squares, from two networks, each giving a tiling by the fourteen
squares of Figure~\ref{fig:tiling}.
\end{proof}

Theorem~\ref{thm:main} follows from the tiling and Proposition~\ref{prop:lower}. Since
$f(1394,1222)\le f(697,611)=14$, the doubled rectangle also needs at most fourteen squares, fewer than
the sixteen previously known.

\section{Checks}

An exhaustive search is only as good as its implementation, so the same program was tested against
the known values. For every rectangle with $m\le n\le 388$ the least number of squares it finds,
taking the best over all divisors of $\gcd(m,n)$, agrees with OEIS A219158 whenever that value is at
most $14$, and finds no tiling when the value exceeds $14$: $65\,231$ rectangles, no disagreement. A
deliberately broken version that discards networks with parallel edges disagrees on $42\,489$ of the
$42\,862$ rectangles with $f\le 12$, so the comparison is sensitive to such errors. The tiling itself is
rechecked by a separate program (\path{verify/verify_tiling.py}) that reads the square sides off a
network, confirms the area, and finds an explicit placement by exhaustive bottom-left filling, using
no network theory.

\begin{remark}
The same enumeration was run up to $16$ squares ($59\,312\,272$ quadrangulations with $19$ vertices).
It finds Pegg's $16$-square tiling of $1394\times1222$ again, which serves as a positive control, and
$241$ further tilings of $697\times611$ by $16$ squares. Within its range, rectangles with sides at most
$1500$ that can be tiled with at most $16$ squares, it finds no rectangle that can be tiled with fewer
squares than a rectangle it is a multiple of.
\end{remark}

\section*{Acknowledgement and data}

The author used an AI tool (Claude, Anthropic) to write and run the enumeration and is responsible for
the content. The tiling was first posted as an answer to the MathOverflow question~\cite{mo} on
29 September 2026. Programs, the OEIS comparison and the plantri 5.8 sources (Apache 2.0 licence) are
in the \path{verify} folder accompanying this note; \path{verify/run.sh} reproduces
Proposition~\ref{prop:lower} in under a minute on a laptop.

\begin{thebibliography}{9}
\bibitem{plantri} G.~Brinkmann and B.~D.~McKay, Fast generation of planar graphs,
\emph{MATCH Commun. Math. Comput. Chem.} \textbf{58} (2007), 323--357.
\bibitem{bsst} R.~L.~Brooks, C.~A.~B.~Smith, A.~H.~Stone and W.~T.~Tutte, The dissection of
rectangles into squares, \emph{Duke Math. J.} \textbf{7} (1940), 312--340.
\bibitem{mo} Wolfgang, Tiling a rectangle with the smallest number of squares, MathOverflow question
116382 (2012), \url{https://mathoverflow.net/q/116382}, accessed 9 October 2026.
\bibitem{pegg} E.~Pegg Jr, answer to~\cite{mo} (9 July 2017), \url{https://mathoverflow.net/a/275053}.
\bibitem{oeis} D.~Radcliffe, Minimum number of integer-sided squares needed to tile an $m\times n$
rectangle, OEIS A219158 and its b-file, \url{https://oeis.org/A219158}, accessed 9 October 2026.
\end{thebibliography}
\end{document}
